% !TeX root = ../main.tex
\section{Introduction}
\label{sec:introduction}

Time-varying conditional quantiles describe how the lower tail, center, and upper tail of a response evolve across an ordered domain. They are useful when changes in variability or asymmetry matter as much as changes in location, as in environmental monitoring, growth assessment, and economic time series. Quantile trend filtering combines the check loss of \citet{koenker1978regression} with an $\ell_1$ penalty on discrete differences, adapting a piecewise polynomial fit to local changes in the signal \citep{tibshirani2014adaptive,madrid2022risk}.

\citet{brantley2020baseline} developed a quantile trend filtering approach for smoothly varying baseline drift, including joint non-crossing estimation, smoothing-parameter selection, and an alternating direction method of multipliers (ADMM) algorithm. Their formulation provides an effective optimization-based benchmark, especially for the smooth-baseline setting that motivates the method. At a selected regularization level, however, the optimizer supplies point estimates; posterior simulation requires a probabilistic model and an inference procedure beyond that optimization step.

We develop a Bayesian extension that retains this computational connection. Under an asymmetric Laplace working likelihood and a global Laplace prior on trend differences, the conditional posterior mode at fixed regularization parameter $\lambda$ and likelihood scale $\sigma$ is the quantile trend filtering solution with penalty $\lambda\sigma$. We infer $\lambda$ and $\sigma$ jointly with the trend instead of selecting the penalty by an external criterion. The posterior mean need not equal the optimizer, but both arise from the same check-loss and difference-penalty structure, and Section~\ref{subsec:individual} states the connection explicitly.

Related methods include Bayesian and calibrated loss-based quantile smoothing. \citet{fasiolo2021fast} developed fast calibrated additive quantile regression. For quantile trend filtering, \citet{onizuka2024fast} combined an asymmetric Laplace likelihood with shrinkage priors and developed both Gibbs sampling and calibrated variational approximations. Their horseshoe prior supports coefficient-specific shrinkage, allowing the amount of regularization to vary across the trend. The supplied Gibbs implementation uses latent-variable and local-scale updates together with a changing Gaussian precision matrix; its matrix operations motivate an alternative computational route for long signals.

Proximal Markov chain Monte Carlo provides such a route for nonsmooth models. The proximal Metropolis-adjusted Langevin algorithm (P-MALA) of \citet{pereyra2016proximal} replaces the nonsmooth negative log-density by its Moreau--Yosida envelope, proposes a Langevin step on the envelope, and corrects with the exact target through a Metropolis--Hastings step; related schemes split the density into a smooth part handled by its gradient and a nonsmooth part handled by a proximal map \citep{durmus2018efficient}. In mean trend filtering, \citet{heng2023bayesian} used proximal Bayesian computation to learn the trend, scale, and regularization parameters. Quantile trend filtering is harder because the check loss is itself nonsmooth, so the proximal map of the full negative log-posterior is not a standard trend filtering problem, and smoothing the check loss to obtain a gradient changes the target that the proposal approximates.

We resolve this with the normal--exponential representation of the asymmetric Laplace distribution \citep{kozumi2011gibbs}. Conditionally on one latent variable per observation, the likelihood is Gaussian in the trend, and the full conditional of the trend is a weighted quadratic plus the $\ell_1$ difference penalty. Its proximal map is a weighted trend filtering problem, which we solve by ADMM with a banded factorization. The resulting Gibbs sampler, which we call Bayesian-PMALA, has conjugate updates for the latent variables, the global penalty $\lambda$ and the scale $\sigma$, and a P-MALA update of the trend built on the proximal map of its whole conditional, with Gaussian noise added after the proximal map so that the Metropolis--Hastings correction is exact. The proposal is preconditioned by the curvature of the conditional, and an ancillarity--sufficiency interweaving step \citep{yu2011center} improves the mixing of $\lambda$. The hierarchy has one global penalty per quantile rather than a vector of local shrinkage scales. One banded factorization serves both the forward and the reverse proximal maps of each move, so for a fixed ADMM iteration budget the work per sweep grows linearly with the sample size. A single default configuration, with difference order $k=2$ and fixed hyperparameters, is used for every signal, noise distribution and quantile level.
% TODO: add the joint non-crossing extension here once Section~\ref{subsec:joint} is rewritten.

Our simulation study focuses on two established benchmarks: the varying-smoothness signal used by \citet{onizuka2024fast} and the sinusoidal smooth signal used by \citet{brantley2020baseline}, under Gaussian, Beta, and mixed-normal noise with the noise parameterization of Onizuka et al. We compare Bayesian-PMALA with the two Onizuka MCMC variants and the ADMM point estimator at quantile levels from $0.05$ to $0.95$, on identical data sets. Each competitor is most accurate on the benchmark that suits it, the horseshoe on the varying-smoothness signal and ADMM on the smooth signal, while Bayesian-PMALA is competitive on both: it attains accuracy comparable to the locally adaptive horseshoe sampler with a simpler hierarchy and banded computations, and accuracy comparable to ADMM while providing posterior inference on the trend, the penalty and the likelihood scale. A third experiment adopts the peak-detection design of \citet{brantley2020baseline}, where the target is a peak-free low-quantile baseline of a contaminated series and estimation is followed by a classification step. Because the likelihood cannot separate a burst from a movement of the trend, this setting uses a stronger prior on the penalty, fixed in advance, with the same sampler; Bayesian-PMALA then recovers the baseline far more accurately than the horseshoe sampler and comes close to ADMM in both baseline recovery and detection.

Bayesian-PMALA thus sits between computationally efficient optimization and a more elaborate locally adaptive Bayesian hierarchy. Relative to ADMM, it adds exact posterior sampling and learns the regularization and likelihood scales instead of selecting a penalty by an information criterion. Relative to the Onizuka Gibbs sampler, it replaces a vector of local shrinkage scales with one global penalty per quantile and replaces dense matrix computations with banded proximal computations.

Section~\ref{sec:methodology} presents the model and Bayesian-PMALA, with computational details in Appendix~\ref{app:pmala}. Section~\ref{sec:simulations} reports the simulation study, with additional tables in Appendix~\ref{app:simulation_results}. Section~\ref{sec:applications} presents the data applications, and Section~\ref{sec:discussion} discusses the findings and limitations.
